\begin{frame}{허브는 더 많이 ``나타나는 단어''}
  \begin{columns}[T]
    \begin{column}{0.42\textwidth}
      \begin{itemize}
        \item 카라테 클럽에서: node 33(관장) 도 17, node 0(감독) 도 16
          --- 두 파벌 리더가 구조적으로 가장 많이 ``나타나는'' 단어.
        \item 임베딩 학습이 자동으로 \textbf{구조적으로 중요한
          노드에 더 많은 ``학습 예산''}을 배정해준다고 볼 수 있다.
      \end{itemize}
      \vskip0.5em
      \begin{block}{주석 --- 13.3과의 연결}
        \small ``영원히 돌아다닌 뒤 머무는 비율''은 댐핑을 빼면
        곧바로 \textbf{PageRank} 그 자체.
      \end{block}
    \end{column}
    \begin{column}{0.55\textwidth}
      \centering
      \includegraphics[width=\linewidth]{karate_embedding_norms.png}
    \end{column}
  \end{columns}

\note{\begin{itemize}
\item 14.2절의 ``\textbf{빈도 높은 단어의 임베딩 벡터 노름이 크다}''
      는 현상과 \textbf{정확히 같은 메커니즘}.
\end{itemize}}
\end{frame}
